\chapter*{CONCLUSION ET PERSPECTIVES}
\markboth{\MakeUppercase{CONCLUSION GÉNÉRALE}}{}
\addcontentsline{toc}{chapter}{CONCLUSION GÉNÉRALE}
\adjustmtc
\thispagestyle{MyStyle}

L'essor du Cloud Computing et des pratiques DevOps a considérablement élargi les
possibilités offertes aux équipes de développement, mais au prix d'une complexité
croissante~: multiplication des technologies à maîtriser, décisions d'architecture
nombreuses et préparation fastidieuse des artefacts de déploiement. C'est dans ce contexte
que s'inscrit ce projet de fin d'études, réalisé au sein de Talan Tunisie Consulting, dont
l'objectif était de concevoir une plateforme capable d'accompagner les ingénieurs Cloud et
DevOps tout au long de ce processus. La solution proposée, \textbf{MultiMindOps}, se
présente comme une plateforme d'assistance fondée sur l'intelligence artificielle~: elle
automatise les tâches répétitives de conception et de préparation de l'infrastructure, tout
en laissant à l'ingénieur la maîtrise des décisions stratégiques.

Au terme de ce travail, la plateforme réunit un ensemble de contributions cohérentes qui,
mises bout à bout, couvrent l'essentiel du cycle d'ingénierie Cloud. À partir d'une simple
analyse automatique d'un projet, MultiMindOps recommande plusieurs architectures cloud
accompagnées d'une estimation des coûts, puis génère de façon cohérente les principaux
artefacts DevOps~--- Docker, Terraform et pipelines GitHub Actions. Sa valeur ne tient pas
seulement à cette automatisation, mais à la fiabilité qu'elle apporte~: une boucle de
remédiation intelligente diagnostique et corrige les erreurs d'exécution, tandis qu'une
logique \textit{Human-in-the-Loop} soumet chaque décision critique à la validation de
l'ingénieur. La sécurité y est traitée comme une exigence de qualité transversale, intégrée
aux artefacts générés plutôt que comme une étape isolée. La plateforme réduit ainsi la
charge cognitive et le temps de mise en place d'une infrastructure cloud, sans jamais
dessaisir l'humain de son rôle de décideur.

\section*{Perspectives}
\addcontentsline{toc}{section}{Perspectives}

Plusieurs pistes réalistes prolongent ce travail. À court terme, l'optimisation du temps
d'exécution des workflows et l'amélioration continue des recommandations grâce aux retours
des utilisateurs renforceraient l'expérience d'usage. À moyen terme, l'extension de la
plateforme à d'autres fournisseurs cloud, tels qu'Azure et Google Cloud Platform,
élargirait son champ d'application au-delà d'AWS. Enfin, l'enrichissement progressif de la
base de connaissances RAG et l'ajout de nouveaux agents spécialisés permettraient de couvrir
davantage de types d'artefacts et de cas d'usage, tout en préservant la cohérence de
l'architecture multi-agents existante.
